\chapter{Objetivos}
\label{chap:objetivos}

\drop{A}{ntes} de abordar el estado del arte y la metodología de desarrollo,
es preciso delimitar con claridad qué pretende lograr este \ac{TFG}.  Este
capítulo establece el objetivo general que guía el desarrollo de \dataqual,
cinco objetivos específicos que cubren las principales dimensiones
funcionales, los objetivos transversales de calidad del software y el
alcance del proyecto con sus limitaciones.


%% =========================================================================
\section{Objetivo general}
\label{sec:objetivo-general}
%% =========================================================================

El objetivo general de este \ac{TFG} consiste en \textbf{diseñar, implementar
y validar una plataforma web interactiva y una \ac{API} \ac{REST}ful que
permitan evaluar de forma automatizada la calidad de los datos empleados en
proyectos de \ac{IA}}, proporcionando métricas cuantitativas, visualizaciones
interactivas y mecanismos de diagnóstico que faciliten la identificación y
corrección de problemas de calidad en los conjuntos de datos.

Este objetivo responde a la necesidad, identificada en la
sección~\ref{sec:motivacion}, de disponer de herramientas que superen las
limitaciones de los enfoques ad~hoc (\emph{scripts} puntuales, consultas
\ac{SQL} manuales o inspecciones en cuadernos Jupyter) y que, al mismo
tiempo, no impongan la barrera técnica que presentan las bibliotecas basadas
exclusivamente en código, como se analizará más adelante en la
sección~\ref{sec:estado-arte}.  \dataqual ofrece una solución
\emph{full-stack}, alineada con los estándares
\ac{ISO}/IEC~25012~\citep{iso25012} e \ac{ISO}/IEC~5259~\citep{iso5259}, que
hace accesible la evaluación de calidad de datos sin necesidad de escribir
código.


%% =========================================================================
\section{Objetivos específicos}
\label{sec:objetivos-especificos}
%% =========================================================================

El objetivo general se descompone en los siguientes cinco objetivos específicos
(\textbf{OE1}--\textbf{OE5}), cada uno de los cuales aborda una dimensión
funcional diferenciada de la plataforma.


%% -------------------------------------------------------------------------
\subsection{OE1.\ Módulo de ingesta y almacenamiento}
\label{sec:oe1}
%% -------------------------------------------------------------------------

Antes de evaluar la calidad de un dataset, la plataforma debe ser capaz de
recibirlo, interpretarlo y custodiarlo de forma fiable.  Este requisito,
aparentemente sencillo, esconde una dificultad práctica recurrente: los
ficheros \ac{CSV} del mundo real presentan combinaciones heterogéneas de
codificaciones, delimitadores y convenciones de formato que los analizadores
estándar no resuelven sin configuración explícita.  Ignorar este problema
implica que una fracción significativa de los datasets reales sería rechazada
en la entrada, invalidando la utilidad de cualquier análisis posterior.

El objetivo consiste en desarrollar un módulo de ingesta y almacenamiento que
resuelva este problema de forma transparente para el usuario, sin requerir
conocimiento técnico sobre el formato del fichero.  Los requisitos que lo
concretan son los siguientes:

\begin{itemize}
  \item \textbf{Soporte \ac{CSV}.}  La plataforma acepta conjuntos de datos
    en formato \ac{CSV}, el formato tabular más extendido en proyectos de
    \ac{IA} y ciencia de datos.

  \item \textbf{Parseo robusto.}  El módulo implementa una estrategia de
    \emph{fallback} que combina múltiples codificaciones (UTF-8, UTF-8-BOM,
    CP1252, Latin-1), separadores (coma, punto y coma, tabulador,
    \emph{pipe}) y motores de lectura, garantizando la ingesta correcta del
    mayor número posible de archivos reales sin intervención del usuario.

  \item \textbf{Almacenamiento seguro.}  Los archivos originales se
    almacenan en un sistema de objetos compatible con S3 (MinIO) para
    garantizar su integridad y disponibilidad; los metadatos (esquema,
    estadísticas descriptivas y conteos) se persisten en PostgreSQL para
    permitir consultas eficientes.

  \item \textbf{Extracción automática de esquema.}  El sistema infiere
    automáticamente los nombres de columna, los tipos de dato y las
    estadísticas descriptivas básicas de cada dataset cargado, eliminando
    la necesidad de configuración manual previa a la evaluación.

  \item \textbf{Versionado de conjuntos de datos.}  El módulo soporta la
    carga de nuevas versiones de un mismo dataset, manteniendo la
    trazabilidad entre versiones mediante una relación padre-hijo y un
    indicador de versión activa.
\end{itemize}


%% -------------------------------------------------------------------------
\subsection{OE2.\ Motor de evaluación con métricas parametrizables}
\label{sec:oe2}
%% -------------------------------------------------------------------------

El núcleo funcional de \dataqual es su motor de evaluación.  Las herramientas
generalistas de calidad de datos ofrecen métricas orientadas a bases de datos
transaccionales (completitud, unicidad, conformidad de formato), pero carecen
de las dimensiones específicas que condicionan el rendimiento de un modelo de
\ac{IA}: el desequilibrio entre clases introduce sesgo sistemático, y la falta
de representatividad de ciertos valores o rangos compromete la capacidad de
generalización.  Este objetivo cubre ambos frentes: las dimensiones estándar
definidas por los organismos normativos y las dimensiones propias del contexto
de \ac{ML}.

El motor debe cubrir las dimensiones de calidad de
\ac{ISO}/IEC~25012~\citep{iso25012}, aplicando las fórmulas de medición
definidas en \ac{ISO}/IEC~25024~\citep{iso25024}, y añadir las métricas
específicas para \ac{IA} de \ac{ISO}/IEC~5259~\citep{iso5259}.  Las \textbf{siete
métricas evaluables} que contribuyen al \emph{Quality Score} son:

\begin{enumerate}
  \item \textbf{Completitud} (\texttt{completeness}): ratio de valores no
    nulos por columna, alineada con la característica de completitud de
    \ac{ISO}/IEC~25012.

  \item \textbf{Exactitud sintáctica} (\texttt{syntactic\_accuracy}):
    conformidad de los valores con patrones de formato definidos mediante
    expresiones regulares, correspondiente a la dimensión de exactitud.

  \item \textbf{Consistencia lógica} (\texttt{logical\_consistency}):
    verificación de reglas condicionales entre columnas (reglas
    \texttt{IF\,--\,THEN}), alineada con la dimensión de consistencia.

  \item \textbf{Actualidad} (\texttt{currentness}): antigüedad del dato más
    reciente respecto a un instante de referencia, correspondiente a la
    dimensión de actualidad (\emph{currentness}).

  \item \textbf{Unicidad} (\texttt{uniqueness}): detección de registros
    duplicados, implementando la sub-dimensión Con-ML-1 de consistencia
    definida en \ac{ISO}/IEC 5259-2.

  \item \textbf{Equilibrio de clases} (\texttt{class\_balance}): evaluación
    del balance de la variable objetivo mediante la entropía de Shannon,
    métrica específica para proyectos de \ac{IA} que permite detectar
    desequilibrios que sesgan los modelos~\citep{ng2021datacentric}.

  \item \textbf{Diversidad del dataset} (\texttt{diversity}): evaluación,
    de forma \emph{opt-in} con expectativas definidas por el usuario, de si
    el dataset cubre los valores y rangos esperados por columna, detectando
    subrepresentación que puede comprometer la capacidad de generalización
    de los modelos~\citep{iso5259}.
\end{enumerate}

Adicionalmente, el motor debe incluir una clase de \textbf{detección de
valores atípicos} (\texttt{outliers}) basada en el \ac{IQR}~\citep{tukey1977}.
Siguiendo la recomendación de \ac{ISO}/IEC~5259~\citep{iso5259} de no tratar
los \emph{outliers} como una dimensión de calidad per~se (su relevancia
depende del dominio de aplicación), esta clase \textbf{no participará en
el \emph{Quality Score}} sino que se utilizará exclusivamente como
herramienta de diagnóstico en el módulo de perfilado de datos (\ac{EDA}).

El motor debe adoptar una arquitectura de \emph{plugins}, de modo que cada
métrica se implemente como una clase independiente que herede de una clase
base abstracta y se registre dinámicamente en un registro centralizado.
Esta arquitectura garantiza la extensibilidad del sistema, permitiendo
incorporar nuevas métricas sin modificar el núcleo del motor.

Adicionalmente, cada métrica debe ser parametrizable: umbrales de aceptación,
columnas objetivo, métodos de detección y pesos relativos (en una escala de
0,0 a 1,0) que determinen su contribución a la puntuación global de calidad.
El sistema debe soportar también plantillas de métricas que permitan
reutilizar configuraciones entre proyectos.


%% -------------------------------------------------------------------------
\subsection{OE3.\ \acs{API} \acs{REST}ful con especificación formal}
\label{sec:oe3}
%% -------------------------------------------------------------------------

Una plataforma de evaluación de calidad de datos no debería ser un sistema
cerrado: su valor aumenta cuando puede integrarse en \emph{pipelines} de
\ac{MLOps}, en procesos de integración continua o en herramientas externas
que consuman sus resultados de forma programática.  Para ello es necesario
que la \ac{API} esté bien estructurada y cuente con una especificación
explícita que permita a terceros integrarla sin necesidad de leer el código
fuente.

El objetivo consiste en diseñar e implementar una \ac{API} \ac{REST}ful que
exponga \emph{endpoints} para la carga de datos, configuración de evaluaciones,
ejecución de análisis y consulta de resultados.  La especificación completa de
rutas, parámetros y formatos de respuesta se recoge en el Anexo~A; este
objetivo define el alcance funcional que dicha especificación debe cubrir.

La \ac{API} se organiza en seis módulos funcionales cohesivos, siguiendo el
patrón \emph{Blueprint} de Flask:

\begin{definitionlist}
  \item[\texttt{auth}] Registro de usuarios, inicio de sesión, refresco de
    \emph{tokens} \ac{JWT} y consulta del perfil del usuario autenticado.

  \item[\texttt{projects}] Operaciones \ac{CRUD} sobre proyectos, incluyendo
    la asignación de conjuntos de datos y la configuración de métricas.

  \item[\texttt{datasets}] Carga de conjuntos de datos, consulta de
    metadatos, gestión de versiones y descarga de archivos.

  \item[\texttt{metrics}] Definición y configuración de métricas de calidad,
    gestión de plantillas y consulta del catálogo de métricas disponibles.

  \item[\texttt{evaluations}] Lanzamiento de evaluaciones, consulta de
    progreso, recuperación de resultados e \emph{issues} detectados.

  \item[\texttt{admin}] Utilidades de administración del sistema, incluyendo
    un \emph{endpoint} de comprobación de salud (\emph{health check}).
\end{definitionlist}

Todos los \emph{endpoints} protegidos requieren autenticación mediante
\ac{JWT}.  El formato de respuesta es uniforme y predecible, con la estructura
\texttt{\{success, data, message\}} y soporte de paginación en los listados.
La convención REST de rutas anidadas añade \emph{blueprints} técnicos
adicionales (\texttt{project\_metrics}, \texttt{project\_datasets},
\texttt{dashboard}) que no amplían el alcance funcional pero sí el número
total de \emph{blueprints} registrados.


%% -------------------------------------------------------------------------
\subsection{OE4.\ Aplicación web con \emph{dashboards} interactivos}
\label{sec:oe4}
%% -------------------------------------------------------------------------

Una \ac{API} potente pero sin interfaz gráfica reproduce exactamente la
barrera de acceso que \dataqual pretende eliminar: limita el uso a perfiles
con conocimientos de programación.  La interfaz web no es un complemento
estético del sistema sino una decisión de diseño central: permite que analistas
de datos, gestores de proyectos o auditores de calidad operen la plataforma
sin escribir una sola línea de código, configurando métricas, interpretando
resultados y tomando decisiones directamente desde el navegador.

El objetivo consiste en desarrollar una aplicación web que cubra el ciclo
completo de trabajo sobre un dataset: gestión de proyectos y versiones,
configuración de evaluaciones, visualización de resultados y exploración
exploratoria de los datos.  Las capacidades concretas que debe ofrecer son:

\begin{itemize}
  \item \textbf{Cuadro de mando principal} con resumen de proyectos,
    conjuntos de datos y distribución de \emph{Quality Gates}.

  \item \textbf{Página de detalle de proyecto} con pestañas diferenciadas
    para conjuntos de datos, configuración de métricas, historial de
    evaluaciones y estado del \emph{Quality Gate}.

  \item \textbf{Visualización de resultados de evaluación} con indicadores de
    puntuación de calidad (\emph{gauge}), pestañas por métrica y listado de
    \emph{issues} detectados.

  \item \textbf{Exploración interactiva de datos} (\ac{EDA}): histogramas,
    diagramas de caja, gráficos de barras, gráficos de dispersión,
    matrices de correlación (Pearson y Spearman) y detección visual de
    valores atípicos con factor \ac{IQR} configurable.

  \item \textbf{Comparación entre versiones de un dataset}: página dedicada
    que muestra las diferencias en número de filas, columnas, puntuación
    de calidad, columnas añadidas o eliminadas e \emph{issues} nuevos,
    resueltos y comunes entre dos versiones.

  \item \textbf{Canvas de linaje de versiones}: árbol visual interactivo
    que representa la genealogía completa de versiones de un dataset,
    permitiendo navegar por la historia de cambios.

  \item \textbf{Gráficos de tendencia} que muestren la evolución de la
    calidad a lo largo de las sucesivas evaluaciones.

  \item \textbf{Indicadores visuales de \emph{Quality Gate}} con códigos de
    color: verde (\texttt{PASSED}), amarillo (\texttt{WARNING}) y rojo
    (\texttt{FAILED}).

  \item \textbf{Diseño adaptable} (\emph{responsive}) con soporte para
    dispositivos de escritorio, tableta y teléfono móvil.
\end{itemize}


%% -------------------------------------------------------------------------
\subsection{OE5.\ Mecanismos de identificación de problemas}
\label{sec:oe5}
%% -------------------------------------------------------------------------

Detectar que un dataset tiene un 15\,\% de valores nulos en una columna es
útil; saber si ese porcentaje ha empeorado respecto a la semana pasada, si el
problema es recurrente o si ya se corrigió una vez, es lo que permite tomar
decisiones informadas.  La mayoría de las herramientas de calidad de datos
producen un informe por ejecución pero no conservan memoria entre ejecuciones:
cada análisis es un snapshot aislado sin conexión con el anterior.

Este objetivo resuelve ese problema introduciendo tres mecanismos que,
combinados, convierten la evaluación puntual en un proceso de seguimiento
continuo:

\begin{enumerate}
  \item \textbf{Sistema de \emph{issues}.}  Cada problema de calidad
    detectado debe registrarse como un \emph{issue} con cuatro niveles de
    severidad (\emph{critical}, \emph{high}, \emph{medium}, \emph{low}),
    calculada dinámicamente en función de la distancia al umbral configurado.
    Cada \emph{issue} debe incluir muestras de los datos afectados (hasta
    cinco filas de ejemplo), con ocultación automática de los valores de
    columnas marcadas como \ac{PII}.

  \item \textbf{\emph{Fingerprinting} determinista.}  Cada \emph{issue} debe
    recibir un identificador único generado mediante \ac{SHA}-256 a partir de
    sus atributos característicos (métrica, columna, tipo de problema),
    permitiendo su identificación unívoca entre ejecuciones independientes.
    Este mecanismo, inspirado en el sistema de rastreo de \emph{issues} de
    SonarQube~\citep{sonarqube}, permite clasificar automáticamente cada
    incidencia como \emph{nueva}, \emph{recurrente} o \emph{corregida}
    respecto a la evaluación anterior.

  \item \textbf{\emph{Quality Gates}.}  Umbrales mínimos configurables que
    determinan si un conjunto de datos es apto para su uso, con tres estados
    posibles: \texttt{PASSED} (score por encima del umbral y sin exceso de
    issues críticos ni nuevos), \texttt{WARNING} (score dentro del margen de
    advertencia configurable por encima del umbral mínimo, o número de issues
    nuevos superior al límite configurado) y \texttt{FAILED} (score por debajo
    del umbral mínimo o issues críticos por encima del máximo permitido).
    Las \emph{Quality Gates} proporcionan una señal clara y automatizable
    que puede integrarse en \emph{pipelines} de
    \ac{MLOps}~\citep{kreuzberger2023} o de integración continua para
    bloquear el uso de conjuntos de datos de calidad insuficiente.
\end{enumerate}


%% =========================================================================
\section{Objetivos transversales}
\label{sec:objetivos-transversales}
%% =========================================================================

Además de los objetivos funcionales descritos, el desarrollo de \dataqual
persigue un conjunto de objetivos transversales de calidad del software que
condicionan las decisiones arquitectónicas y tecnológicas del proyecto:

\begin{definitionlist}
  \item[Seguridad]  Autenticación basada en \ac{JWT} con \emph{tokens} de
    acceso y refresco, \emph{hashing} de contraseñas mediante \ac{PBKDF2},
    limitación de tasa de peticiones (\emph{rate limiting}), validación de
    entradas en los puntos críticos mediante esquemas Marshmallow, protección
    contra inyección de código en las reglas de consistencia lógica y
    enmascaramiento automático de columnas marcadas como \ac{PII}.

  \item[Escalabilidad]  Procesamiento asíncrono de evaluaciones mediante
    Celery y Redis como \emph{broker} de mensajes, lo que permite escalar
    horizontalmente el número de \emph{workers} en función de la carga de
    trabajo, así como separación de responsabilidades entre servicios.

  \item[Mantenibilidad]  Arquitectura modular basada en \emph{blueprints},
    servicios y modelos, con un sistema de métricas extensible por
    \emph{plugins} que permite incorporar nuevas métricas con un esfuerzo
    mínimo.  Migraciones de base de datos versionadas y separación estricta
    entre \emph{frontend} y \emph{backend}.

  \item[Usabilidad]  Interfaz de usuario basada en Material Design con
    formularios dotados de validación en tiempo real, notificaciones
    \emph{toast}, migas de pan (\emph{breadcrumbs}) de navegación y estados
    de carga y error claramente diferenciados.  La plataforma soporta
    español e inglés con cambio de idioma en tiempo real, sin necesidad de
    recargar la aplicación.

  \item[Portabilidad]  Despliegue del \emph{stack} completo mediante Docker
    Compose, con configuración por variables de entorno y sin dependencias de
    servicios en la nube propietarios.  Los ocho servicios que componen la
    plataforma (\emph{frontend}, \emph{backend}, PostgreSQL, MinIO, Redis,
    Celery \emph{worker}, Celery Beat y Flower) se orquestan conjuntamente,
    lo que garantiza la reproducibilidad del entorno y que los datos no
    salgan de la infraestructura propia.
\end{definitionlist}


%% =========================================================================
\section{Alcance y limitaciones}
\label{sec:alcance-objetivos}
%% =========================================================================

El alcance del proyecto abarca el diseño, la implementación y la validación
de \dataqual como plataforma \emph{full-stack}, cubriendo los cinco objetivos
específicos descritos en la sección~\ref{sec:objetivos-especificos}: el módulo
de ingesta con parseo robusto y versionado, el motor de evaluación con siete
métricas parametrizables, la \ac{API} \ac{REST}ful con los módulos funcionales
planificados, la aplicación web con cuadros de mando interactivos, y los
mecanismos de identificación y seguimiento de problemas.  El sistema se
despliega íntegramente mediante Docker Compose, se valida con conjuntos de
datos reales y se documenta en los anexos de la presente memoria
(especificación de la \ac{API} y catálogo de métricas).  La pila tecnológica
queda fijada por las decisiones de arquitectura adoptadas: Python con Flask
y PostgreSQL en el \emph{backend}, React con Next.js y TypeScript en el
\emph{frontend}, Celery con Redis para el procesamiento asíncrono, y
alineación explícita con los estándares \ac{ISO}/IEC~25012~\citep{iso25012}
e \ac{ISO}/IEC~5259~\citep{iso5259}.

Fuera de este alcance quedan tres aspectos que, si bien son relevantes en
un contexto de producción real, no forman parte del presente trabajo.

El primero es el \textbf{procesamiento de grandes volúmenes de datos}.  El
motor de evaluación carga el dataset completo en memoria mediante Pandas, lo
que limita el tamaño práctico de los ficheros a lo que la memoria RAM disponible
permita procesar sin degradación.  La evaluación distribuida sobre datasets
masivos mediante Apache Spark o Dask requiere una arquitectura diferente y
queda como línea de trabajo futuro.

El segundo es la \textbf{corrección automática de datos}.  La plataforma
diagnostica y clasifica los problemas de calidad, pero la decisión y
ejecución de las acciones correctoras (limpieza, imputación, deduplicación)
corresponde al usuario.  Incorporar sugerencias de corrección automatizadas
es una de las líneas de trabajo futuro identificadas en el
capítulo~\ref{chap:conclusiones}.

El tercero es la \textbf{integración directa con fuentes de datos externas}.
La arquitectura de almacenamiento basada en MinIO es compatible con el
protocolo S3, pero no se desarrollan conectores específicos para bases de
datos (MySQL, BigQuery, Snowflake) ni para servicios de almacenamiento en la nube
(AWS S3, Azure Blob, Google Cloud Storage).  El usuario interactúa con la
plataforma a través de ficheros \ac{CSV} cargados manualmente.


% Local Variables:
%  coding: utf-8
%  mode: latex
%  mode: flyspell
%  ispell-local-dictionary: "castellano8"
% End:
